% =============================================================================
% CompetitiveExam/CS301/06-linked-lists-singly-answers.tex
% Answer key for 06-linked-lists-singly-questions.tex.
%
% Real GATE PYQs (Q1--Q6) are sourced from the local GATEOverflow Volume 2
% PYQ book (CompetitiveExam/Books/index.md, Tier-1 availability: Linked List
% cluster PDF pp.232-238), which authenticates each question (year + question
% number printed on the page). The book's own answer keys are collapsed and
% not extractable, so each real-PYQ answer was independently established in
% the build session (2026-09-26): GATEOverflow online question pages,
% ExamSIDE year-wise/topic listings, solutionsadda's GATE-1994 listing, the
% edurev GATE CS 1994/2002 papers, the gatesolutions blog (2002 Q1.5),
% electrical4u's GATE CS 2003/2008/2010 solutions, Knowledge Gate's SLL PYQ
% page (2003 Q90), the EduRev 2008 Q62 discussion, the GFG Linked-Lists PYQ
% quiz (2010 Q36), compsciedu's MCQ pages, the Collegedunia official
% GATE-2022 paper+solutions PDF, and Garg's Academy GATE-2022 DS solutions.
% Original questions (Q7--Q8) have answers derived directly from the Week 6
% deck algorithms, confirmed by execution (2026-09-26).
% =============================================================================

\documentclass[11pt]{article}
\usepackage[margin=1in]{geometry}
% =============================================================================
% Preambles/header.tex — shared LaTeX/Beamer preamble for this project
%
% Use from a Beamer lecture by prepending:
%     \input{header}
% and compiling with TEXINPUTS=../Preambles:$TEXINPUTS (the /compile-latex
% skill does this automatically).
%
% PALETTE CONTRACT. Color names defined here MUST match the names in
% Quarto/theme-template.scss so Beamer and Quarto renderings use the same
% palette. The scripts/check-palette-sync.sh script enforces this.
%
% To customize: edit the HEX values below AND the matching $variable in
% Quarto/theme-template.scss, then run scripts/check-palette-sync.sh.
% =============================================================================

% -----------------------------------------------------------------------------
% Engine detection + encoding
%
% Our /compile-latex skill uses XeLaTeX, where inputenc is unnecessary and
% can produce encoding warnings. Only load inputenc under pdfTeX.
% -----------------------------------------------------------------------------
\usepackage{iftex}
\ifPDFTeX
  \usepackage[utf8]{inputenc}
\fi

\usepackage{amsmath, amssymb, amsthm}
\usepackage{booktabs}
\usepackage{graphicx}

% -----------------------------------------------------------------------------
% PALETTE — mirrors Quarto/theme-template.scss variable names.
% Load xcolor and define colors BEFORE anything that references them
% (notably hyperref, hypersetup, and the Beamer color assignments).
% -----------------------------------------------------------------------------
\usepackage{xcolor}

\definecolor{primary-blue}{HTML}{012169}      % headings, accents
\definecolor{primary-gold}{HTML}{B9975B}      % emphasis, borders
\definecolor{highlight-yellow}{HTML}{F2A900}  % markers, alerts
\definecolor{light-bg}{HTML}{E8EDF5}          % title / section backgrounds
\definecolor{jet}{HTML}{1A1A1A}               % body text

% Semantic colors (used in diagrams and annotations)
\definecolor{positive}{HTML}{15803D}          % good / observed / identified
\definecolor{negative}{HTML}{B91C1C}          % bad / problematic / bias
\definecolor{neutral}{HTML}{525252}           % reference / context

% Highlight accents (match .hi-* classes in the SCSS)
\definecolor{hi-slate}{HTML}{314F4F}
\definecolor{hi-green}{HTML}{2E8B57}
\definecolor{hi-red}{HTML}{C41E3A}

% Semantic aliases so skill templates and snippets can write
% \textcolor{accent}{...} without hard-coding "primary-blue".
\colorlet{accent}{primary-blue}
\colorlet{accent2}{primary-gold}

% -----------------------------------------------------------------------------
% Hyperref — loaded AFTER xcolor + palette so link colors resolve.
% -----------------------------------------------------------------------------
\usepackage{hyperref}
\hypersetup{colorlinks=true, linkcolor=primary-blue, urlcolor=primary-blue,
            citecolor=primary-blue, pdfborder={0 0 0}}

% -----------------------------------------------------------------------------
% Course code — set once per deck (after \input{header}, before \begin{document})
% via \coursecode{CS401}. Shown in the footer on every slide so a deck's course
% is identifiable without opening the title page. Empty by default.
% -----------------------------------------------------------------------------
\makeatletter
\newcommand{\coursecode}[1]{\gdef\@coursecodetext{#1}}
\gdef\@coursecodetext{}
\makeatother

% -----------------------------------------------------------------------------
% Beamer theme assignments (only apply under Beamer).
%
% We detect Beamer via \@ifclassloaded{beamer}, which is the documented,
% non-internal way. This must be wrapped in \makeatletter/\makeatother
% because \@ifclassloaded contains an @-catcode character.
% -----------------------------------------------------------------------------
\makeatletter
\@ifclassloaded{beamer}{%
  \usefonttheme{professionalfonts}
  \setbeamercolor{structure}{fg=primary-blue}
  \setbeamercolor{frametitle}{fg=primary-blue}
  \setbeamercolor{title}{fg=primary-blue}
  \setbeamercolor{subtitle}{fg=primary-gold}
  \setbeamercolor{author}{fg=primary-blue}
  \setbeamercolor{institute}{fg=neutral}
  \setbeamercolor{date}{fg=primary-gold}
  \setbeamercolor{itemize item}{fg=primary-blue}
  \setbeamercolor{itemize subitem}{fg=primary-gold}
  \setbeamercolor{alerted text}{fg=primary-gold}
  \setbeamercolor{block title}{fg=white, bg=primary-blue}
  \setbeamercolor{block body}{bg=light-bg}
  % Semantic block variants: block=definition/concept (blue), exampleblock=
  % worked example (green/positive), alertblock=key takeaway or caution (gold).
  % Keep to <=2 colored boxes per slide across ALL three variants (INV-7).
  \setbeamercolor{block title example}{fg=white, bg=positive}
  \setbeamercolor{block body example}{bg=light-bg}
  \setbeamercolor{block title alerted}{fg=white, bg=primary-gold}
  \setbeamercolor{block body alerted}{bg=light-bg}

  % Minimal footer: course code (left) + page X / N (right), no navigation clutter
  \setbeamertemplate{navigation symbols}{}
  \setbeamertemplate{footline}{%
    \color{neutral}\footnotesize\hspace{0.7em}\@coursecodetext\hfill
    \insertframenumber/\inserttotalframenumber\hspace{0.7em}\vspace{0.3em}}
}{}
\makeatother

% -----------------------------------------------------------------------------
% TikZ setup — libraries the snippet gallery uses
% -----------------------------------------------------------------------------
\usepackage{tikz}
\usetikzlibrary{arrows.meta, positioning, calc, decorations.pathreplacing,
                fit, shapes.geometric, backgrounds}

% Shared TikZ styles — snippets and hand-written diagrams can pick these up
% via `\begin{tikzpicture}[every node/.style={...}]` or by naming specific
% styles (e.g., `\node[dag-node] (X) at (0,0) {X};`).
\tikzset{
  % Node styles for causal DAGs and similar
  dag-node/.style = {
    draw=primary-blue, fill=light-bg, rounded corners=2pt,
    minimum width=1.2cm, minimum height=0.9cm, align=center, font=\small
  },
  decision-node/.style = {
    draw=primary-gold, fill=white, diamond, aspect=1.8,
    minimum width=1.8cm, minimum height=1.2cm, align=center, font=\small,
    inner sep=1pt
  },
  % Edge styles — observed vs counterfactual
  observed-edge/.style = {-{Latex[length=2.5mm]}, thick, draw=jet},
  counterfactual-edge/.style = {-{Latex[length=2.5mm]}, thick, draw=neutral, dashed},
  confound-edge/.style = {-{Latex[length=2.5mm]}, thick, draw=negative, dashed},
  % Data-point styles for plots
  observed-dot/.style = {fill=primary-blue, draw=primary-blue, circle, inner sep=1.2pt},
  counterfactual-dot/.style = {fill=white, draw=neutral, circle, inner sep=1.2pt}
}

% -----------------------------------------------------------------------------
% Convenience macros
% -----------------------------------------------------------------------------
\newcommand{\muted}[1]{\textcolor{neutral}{#1}}
\newcommand{\key}[1]{\textcolor{primary-gold}{\textbf{#1}}}
\newcommand{\good}[1]{\textcolor{positive}{#1}}
\newcommand{\bad}[1]{\textcolor{negative}{#1}}

% Standout frame macro: full-bleed dark background for section transitions.
% Beamer-only (uses \begin{frame}/\setbeamercolor/\usebeamerfont) -- do not
% call from an article-class file (e.g. Notes/<CODE>/*-notes.tex). \muted,
% \key, \good, \bad, and \coursecode above are plain xcolor/gdef and are
% safe to use from either class; verified when Notes/article-class support
% was added.
% Expand: \transitionslide{Your transition title}
\newcommand{\transitionslide}[1]{%
  \begingroup
  \setbeamercolor{background canvas}{bg=primary-blue}
  \begin{frame}[plain]
    \centering
    \vfill
    {\usebeamerfont{title}\color{white}\Large #1\par}
    \vfill
  \end{frame}
  \endgroup
}

% Section divider frame (Beamer-only), an alternative to \transitionslide with a
% darker two-line style: near-black (jet) background, a small white label line
% above a larger gold title, and a thin gold rule along the bottom edge.
% Expand: \sectiondivider{Section 2}{Pointers}
\newcommand{\sectiondivider}[2]{%
  \begingroup
  \setbeamercolor{background canvas}{bg=jet}
  \begin{frame}[plain]
    \vspace*{3.0cm}
    {\color{white}\ttfamily\large #1\par}
    \vspace{0.5cm}
    {\color{primary-gold}\LARGE #2\par}
    \begin{tikzpicture}[remember picture, overlay]
      \draw[primary-gold, line width=2.5pt]
        ($(current page.south west)+(0,0.1)$) -- ($(current page.south east)+(0,0.1)$);
    \end{tikzpicture}
  \end{frame}
  \endgroup
}

% -----------------------------------------------------------------------------
% Channel watermark — "Concepts That Click" (YouTube).
%
% Article-class files (CompetitiveExam Q/A sets, Notes, Assignments) get a
% light diagonal draft watermark on every page plus a centered footer naming
% the channel. Beamer decks get a subtle TikZ background watermark instead
% (the footline already carries the course code). Centralized here so every
% MCQ PDF is branded without per-file edits.
% -----------------------------------------------------------------------------
\makeatletter
\@ifclassloaded{article}{%
  \usepackage{draftwatermark}
  \SetWatermarkText{Concepts That Click}
  \SetWatermarkScale{1}
  \SetWatermarkFontSize{2cm}
  \SetWatermarkLightness{0.86}
  \SetWatermarkAngle{45}
  \usepackage{fancyhdr}
  \pagestyle{fancy}
  \fancyhf{}
  \fancyfoot[C]{\footnotesize\textcolor{neutral}{Concepts That Click~|~YouTube: \href{https://www.youtube.com/@concepts.thatclick}{@concepts.thatclick}}}
  \renewcommand{\headrulewidth}{0pt}
  \renewcommand{\footrulewidth}{0pt}
  \fancypagestyle{plain}{%
    \fancyhf{}%
    \fancyfoot[C]{\footnotesize\textcolor{neutral}{Concepts That Click~|~YouTube: \href{https://www.youtube.com/@concepts.thatclick}{@concepts.thatclick}}}%
    \renewcommand{\headrulewidth}{0pt}%
    \renewcommand{\footrulewidth}{0pt}%
  }
}{}
\@ifclassloaded{beamer}{%
  \setbeamertemplate{background}{%
    \begin{tikzpicture}[remember picture, overlay]
      \node[rotate=45, opacity=0.055, align=center]
        at (current page.center)
        {\Huge\textcolor{neutral}{Concepts That Click}};
    \end{tikzpicture}%
  }
}{}
\makeatother 
\coursecode{CS301}
\renewcommand{\thesection}{6.\arabic{section}}

\title{Competitive Exam Practice 6 --- Answer Key\\\large Singly Linked Lists}
\author{Auqib Hamid Lone\\[0.3em]
  \emph{Department of Computer Science \& Engineering}, \emph{GCET Ganderbal}}
\date{\today}

\begin{document}
\maketitle

\begin{enumerate}
\item[\textbf{Q1}] \textbf{(B) Binary search.} Binary search needs $O(1)$ random access to the middle element, but reaching the middle of a linked list costs a linear walk --- so the halving step is $O(n)$, not $O(1)$, and the $O(\log n)$ bill collapses. Insertion sort (splice without shifting), radix sort (sequential bucket appends), and polynomial manipulation (one node per nonzero term) all suit lists. Confirmed on the GATEOverflow page, solutionsadda's GATE-1994 listing, the edurev GATE CS 1994 paper, and ExamSIDE's GATE-1994 listing (2026-09-26).

\item[\textbf{Q2}] \textbf{(D) $n$.} A singly linked list has no indexing, so only linear search applies: the miss (or hit at the last node) compares against all $n$ elements. $n/2$ is the \emph{average}-case figure, and the $\log$ options wrongly assume binary search, which needs $O(1)$ mid-access. Confirmed on the GATEOverflow page (answer D), the gatesolutions blog (answer (d) $n$), the edurev GATE CS 2002 paper (Q1.5), and ExamSIDE's GATE-2002 listing (2026-09-26).

\item[\textbf{Q3}] \textbf{(B) The elements are sorted in non-decreasing order.} The recursion demands \texttt{p->data <= p->next->data} at \emph{every} adjacent pair, which is exactly non-decreasing order (equal neighbours allowed). (A) is contained in (B) but not equivalent --- \texttt{f} also returns $1$ on longer sorted lists, so ``if and only if'' rules (A) out; (C) flips the comparison; (D) is neither necessary nor sufficient. Confirmed on the GATEOverflow page (answer B), ExamSIDE's GATE-2003 listing, electrical4u's GATE CS 2003 solutions (answer B), and Knowledge Gate's SLL PYQ page (2026-09-26).

\item[\textbf{Q4}] \textbf{(B) $2, 1, 4, 3, 6, 5, 7$.} Each pass swaps one adjacent pair by \emph{value} and then skips both nodes (\texttt{p = q->next}), so pairs $(1,2)$, $(3,4)$, $(5,6)$ are swapped and the odd node $7$ has no partner and stays put. (A) would need zero swaps; (C) shifts the pairing by one; (D) rotates instead of swapping. Confirmed on the GATEOverflow page (answer B), ExamSIDE's GATE-2008 listing, electrical4u's GATE CS 2008 solutions (answer B), and the EduRev discussion (2026-09-26).

\item[\textbf{Q5}] \textbf{(D) \texttt{q->next = NULL; p->next = head; head = p;}.} On exit \texttt{p} is the last node and \texttt{q} its predecessor: detach the old tail link first (\texttt{q->next = NULL}), point the last node at the old head, then swing \texttt{head} over --- in exactly this order. (A) orphans the whole rest of the list (\texttt{q} is dropped, tail never detached); (B) sets \texttt{head = p} before linking so \texttt{p->next = head} is a self-loop; (C) mislinks \texttt{p->next} to \texttt{q}. Confirmed on the GATEOverflow page (answer D), the GFG PYQ quiz, compsciedu's MCQ page (answer (d)), electrical4u's GATE CS 2010 solutions (answer D), and ExamSIDE's GATE-2010 listing (2026-09-26).

\item[\textbf{Q6}] \textbf{(A) The best algorithm takes $\Theta(n)$ time in the worst case.} The deck's three-pointer loop visits each node exactly once and flips one link per node --- $n$ flips, $\Theta(n)$ worst case, with only three extra pointers ($O(1)$ space). No $o(n)$ algorithm can touch all $n$ links, ruling out (B)/(C); (D) is refuted by the very loop above. Confirmed on the GATEOverflow page (answer A), the official Collegedunia GATE-2022 paper+solutions PDF (correct answer A), Garg's Academy solutions (answer A), ExamSIDE's GATE-2022 listing, and GFG's quiz page (2026-09-26).

\item[\textbf{Q7}] \textbf{(A) \texttt{newNode->next = head;} then \texttt{head = newNode;} giving \texttt{head} $\to$ 4 $\to$ 11 $\to$ 23 $\to$ \texttt{NULL}.} The new box must point at the old chain \emph{before} \texttt{head} moves; reversing the order ((B)/(D)) overwrites \texttt{head}'s old value with no copy anywhere and orphans $11 \to 23$. (C) keeps the safe order but misstates the result (front-insert never appends). Confirmed by execution (2026-09-26).

\item[\textbf{Q8}] \textbf{(A) The list is \texttt{head} $\to$ 8 $\to$ 29 $\to$ \texttt{NULL}; the bypass is $O(1)$ worst case once the predecessor is in hand, and the overall delete is $O(n)$ worst case because of the search.} Node $8$'s \texttt{next} is rewritten to skip $13$ straight to $29$, then $13$'s box is freed. (B) leaves the list unchanged; (C) deletes the wrong node and claims the search is free; (D) gets the list right but drops the $O(n)$ search bill the deck's price list charges. Confirmed by execution (2026-09-26).
\end{enumerate}

\end{document}
